\begin{lemma}[Transporting an exact Tashkinov augmentation route]
Let \(S,B,X\) be ordered tree-sequences for the same distinguished edge,
and let \(\varphi\) and \(\sigma\) be two partial edge colorings on the
same finite loopless multigraph.  Suppose that
\[
 S\longrightarrow X
\]
is a \(\varphi\)-\noderef{TashkinovAugmentationSequence}, that
\(B\preccurlyeq X\) is an exact
\noderef{TreeSequenceInitialSegment} with
\(|X|=|B|+1\), and that the following two pointwise conditions hold.
\begin{enumerate}
\item Every exact edge of \(X\) at a position \(t\) with
\(|S|\leq t<|X|\) retains its color:
\(\sigma(X.\mathit{edge}(t))=\varphi(X.\mathit{edge}(t))\).
\item At every vertex \(v\in\noderef{TreeSequenceVertexSet}(B)\),
\[
 \noderef{MissingColors}(G,\sigma,v)
   =\noderef{MissingColors}(G,\varphi,v).
\]
\end{enumerate}
Then the identical ordered chain is a
\(\sigma\)-\noderef{TashkinovAugmentationSequence} from \(S\) to \(X\).
\end{lemma}
\begin{proof}
We first record the exact-prefix bookkeeping used at every step.  Exact
initial segments compose: if \(A\preccurlyeq C\preccurlyeq D\), compose
the two length inequalities, and at each retained vertex or edge index
compose the two entrywise equalities in
\noderef{TreeSequenceInitialSegment}.  Hence induction on the
reflexive-transitive chain in
\noderef{TashkinovAugmentationSequence} shows that the source and every
intermediate member of a route are exact initial segments of its target.

Let \(A\to C\) be a one-step
\noderef{TashkinovAugmentation} occurring in the given route, and retain
the remaining suffix \(C\to X\).  The one-step length equation gives
\(|C|=|A|+1\), while the suffix gives \(C\preccurlyeq X\).  Thus
\(|A|<|X|=|B|+1\), so \(|A|\leq|B|\).  Both \(A\) and \(B\) are exact
initial segments of \(X\).  For each index of \(A\), compare its entry
with the entry of \(X\), then with the corresponding entry of \(B\).
This proves the exact inclusion
\[
                         A\preccurlyeq B.                    \tag{1}
\]
In particular, \(B\) contains every augmentation predecessor, including
all of its individual displayed vertices.

Unpack the one-step augmentation.  It supplies a color \(c\), the exact
new edge
\[
 g=C.\mathit{edge}(|A|),\qquad \varphi(g)=c,                 \tag{2}
\]
and a witness to
\(c\in\noderef{TreeMissingColors}(\varphi,A)\).  By the definition of
tree-missing colors, choose a particular
\(v\in\noderef{TreeSequenceVertexSet}(A)\) at which \(c\) is missing.
The exact inclusion (1) puts the same vertex identity in
\(\noderef{TreeSequenceVertexSet}(B)\).  Pointwise missing-set equality
on \(B\) therefore says that \(c\) is missing at this same \(v\) under
\(\sigma\).  Reusing the original membership of \(v\) in the predecessor
gives
\[
                   c\in\noderef{TreeMissingColors}(\sigma,A). \tag{3}
\]

The edge identity in (2) occurs at position \(|A|\) in \(X\), because
\(C\preccurlyeq X\) retains every exact edge.  Along the route the lengths
never decrease, so \(|S|\leq|A|\).  The fixed-edge premise at precisely
this position gives \(\sigma(g)=\varphi(g)=c\).  Together with (3), the
unchanged one-step length and exact-prefix clauses prove
\[
                  \noderef{TashkinovAugmentation}(\sigma,A,C). \tag{4}
\]

Finally perform head induction on the original reflexive-transitive
chain.  Maintain the numerical invariant \(|S|\leq|A|\) at the current
source; (4) transports the head step, and its unit length increase gives
the invariant for the next source.  The reflexive case is reflexive under
\(\sigma\), and the induction step prepends (4) to the transported suffix.
This constructs the claimed \(\sigma\)-augmentation sequence with the
same intermediate tree-sequences and the same individual added edges.
\end{proof}
